\section{Introduction}
\label{sec:intro}
\begin{table*}
  \centering
  \resizebox{\linewidth}{!}{ 
  \begin{tabular}{l|cccc|ccc|cc|c}
    \toprule
    Dataset & Annotation & \#Videos & Avg. video len & \#Avg. event & Vision & Audio & Subtitle & Captions & Timestamps & A-V Correlations  \\
    \midrule
    
    InternVid~\cite{wang2023internvid} & G & 234M & 11.7s & 1 & \checkmark & \ding{55}& \ding{55}& V & - & \ding{55}   \\
    Panda-70M~\cite{chen2024panda} & G & 70.8M & 8.5s & 1 & \checkmark & \ding{55}& \ding{55}& V & - & \ding{55} \\
    AudioCaps~\cite{kim2019audiocaps} & M & 51.3K & 10s & 1 & \ding{55} & \checkmark & \ding{55} & A & - & \ding{55} \\
    WavCaps~\cite{mei2024wavcaps} & G & 403K & 67.6s & 1 & \ding{55} & \checkmark & \ding{55} & A & - & \ding{55}\\
    ACAV~\cite{lee2021acav100m} & G & 100M & 10s & 1 & \checkmark & \checkmark & \ding{55} & - & - & \ding{55} \\
    VALOR~\cite{liu2024valor} & M & 1.18M & 10s & 1 & \checkmark & \checkmark & \ding{55} & VA & - & \ding{55}\\
    VAST~\cite{chen2023vast} & G & 27M & 5$\sim$30s & 1 & \checkmark & \checkmark & \checkmark & VAS & - & \ding{55} \\
    \midrule
    AVEL~\cite{tian2018audio} & M & 4,143 & 10s & 1 & \checkmark & \checkmark & \ding{55} & - & VA & \ding{55}\\
    UnAV-100~\cite{geng2023dense} & M & 10,790 & 42.1s & 2.8 & \checkmark & \checkmark & \ding{55} & - & VA & \ding{55} \\
    ActivityNet Caps~\cite{krishna2017dense}& M & 20K & 180s & 3.7 &  \checkmark & \ding{55} & \ding{55} & V & V & \ding{55}\\
    Charades-STA~\cite{gao2017tall} & M & 10K & 30s & 1.6 &  \checkmark & \ding{55} & \ding{55} & V & V & \ding{55}\\
    \midrule
    \textbf{LongVALE (Ours)} & G+M & 8,411 & 235s & 12.6 & \checkmark & \checkmark & \checkmark & VAS & VAS  & \checkmark\\
    \bottomrule
  \end{tabular}}
  \vspace{-2mm}
  \caption{Comparison of LongVALE with previous related benchmarks. G: generated. M: manual. V: visual. A: audio. S: speech. }
  \label{tab:dataset_compar}
  \vspace{-5mm}
\end{table*}

As the volume of videos on social media platforms grows exponentially, video understanding~\cite{krishna2017dense,gao2017tall,wang2023internvid,zhang2023video} has emerged as a vital research area in artificial intelligence.
When watching videos, such as daily vlogs or tutorials lasting several minutes, viewers need to integrate visual and auditory information and associate multiple events to fully comprehend the content. 
An ideal intelligent video agent should imitate it, capable of both cross-modal reasoning and fine-grained temporal understanding. 
However, current research is limited to coarse-grained tasks (\eg, video retrieval/captioning~\cite{wang2023internvid,liu2024valor}) or visual-only fine-grained tasks (\eg, temporal grounding/dense captioning~\cite{krishna2017dense, gao2017tall}), remaining far from enough to achieve both the capabilities.

A significant barrier to this advancement is the absence of a high-quality video dataset with omni-modality (vision, audio, and speech) and fine-grained temporal annotations. 
As seen in Tab.~\ref{tab:dataset_compar}, current benchmarks either contain only global captions for short video/audio clips~\cite{wang2023internvid, chen2024panda, kim2019audiocaps,mei2024wavcaps,liu2024valor,chen2023vast}, offer visual-only multi-event annotations~\cite{krishna2017dense,gao2017tall}, or possess multi-modal events but lack detailed captions~\cite{tian2018audio,geng2023dense}.
Moreover, creating such a dataset
poses significant challenges, as identifying temporal boundaries and producing detailed event captions by integrating information from various modalities are difficult and time-consuming, even for human annotators.  

In this work, we propose an efficient and scalable annotation pipeline, capable of generating temporal boundaries and detailed captions for omni-modal events (\ie, integrating vision, audio, and speech) within arbitrary multi-modal long videos.
Our pipeline includes three distinct aspects: 
1) \textit{High-quality video filtering for rich audio-visual semantics and temporal dynamics.}
The filtered videos showcase rich dynamic visual scenes paired with diverse audio types, \eg, instruments playing, people laughing, and tools whirring as in Fig.~\ref{fig:fig1}, contrasting with the dense narration in prior datasets~\cite{chen2024panda, chen2023vast}. 
2) \textit{Omni-modal event boundary detection for semantic coherence in both visual and audio scenes.} Unlike previous works~\cite{krishna2017dense,gao2017tall} that only identify visual event boundaries, we determine omni-modal event boundaries utilizing both visual and audio cues. This prevents audio scenes from being cut off, avoiding the loss of critical information.
3) \textit{Omni-modal event captioning emphasizing audio-visual correlation reasoning.} Instead of simple concatenation~\cite{chen2023vast,liu2024valor}, we fully integrate modality-specific information (vision, audio, speech) and explicitly reason about their correlations to enhance cross-modal understanding.
For example, in Fig.~\ref{fig:fig1}, the visible man using the tool with loud whirring reflects audio-visual \textit{synchronicity}, and the woman's speech crucially \textit{complements} the visual scene. 

Based on the pipeline, we construct LongVALE, the first-ever benchmark for omni-modality fine-grained video understanding. It comprises 8.4K long videos containing 105K omni-modal events annotated with high-quality temporal boundaries and correlation-aware captions. 
Notably, it features a longer average video length (235 seconds) and more events (12.6 per video) compared to existing benchmarks, along with its unique omni-modal event boundaries and captions with audio-visual correlations (seen in Tab.~\ref{tab:dataset_compar}).


Building upon LongVALE, we present LongVALE-LLM, a multi-modal video LLM, capable of both cross-modal reasoning and fine-grained temporal understanding. Further, we extend three fine-grained video tasks (\ie, temporal video grounding, dense video captioning, and segment captioning) from a vision-oriented to a novel omni-modal setting.
During training, our dataset can serve as a highly valuable data source for both event boundary perception tuning~\cite{huang2024vtimellm} and instruction tuning.
Our experiments show that the model trained on our LongVALE dataset significantly outperforms existing video LLMs across all three tasks.
Moreover, we find that our model can surprisingly achieve superior performance on general audio-visual question answering (AVQA) tasks~\cite{alamri2019audio,li2022learning} in a zero-shot manner, even with significantly less data compared to other LLM-based methods trained on million-scale data. 
It highlights the effectiveness and great promise of our dataset in diving forward comprehensive multi-modal video understanding.
\begin{figure*}[t]
  \centering
  \setlength{\abovecaptionskip}{1.0mm}
   \includegraphics[width=1.0\linewidth]{figures/new/fig2.jpg}
   \caption{The pipeline for high-quality omni-modality fine-grained data generation. 
   It starts by detecting visual and audio event boundaries based on their distinct properties. 
   Next, we generate detailed captions for each video and audio event enhanced by keyframe and speech captions. 
   We then determine omni-modal event boundaries by maintaining the semantic integrity of single-modal events. Finally, omni-modal event captions are generated by audio-visual correlation reasoning, followed by manual refinement to ensure data's high quality.}
   \label{fig:dataset_pipeline}
     \vspace{-4mm}
\end{figure*}
Our contributions can be summarized as follows: 
\begin{itemize}
    \item We propose a novel scalable pipeline enabling the automatic generation of high-quality omni-modality fine-grained annotations for multi-modal long videos, significantly reducing manual annotation costs. 
    \item We introduce LongVALE, the first-ever benchmark providing omni-modal event temporal boundaries and cross-modal correlation-aware captions for 105K omni-modal events within 8.4K high-quality multi-modal long videos.
    \item We demonstrate that our LongVALE-trained model excels in both cross-modal reasoning and fine-grained temporal understanding, significantly outperforming existing video LLMs across all three omni-modal tasks and even achieving superior zero-shot results on general AVQA.
    
\end{itemize}
